% =============================================================================
% 07_longitudinal_registration.tex -- Longitudinal Registration (~6-8 frames)
% =============================================================================

\section{Longitudinal Registration}

% ---- Frame 1 ----
\begin{frame}{The Longitudinal Registration Problem}

  \textbf{Goal:} Align pre-op and post-op T1w volumes to quantify
  tissue displacement between sessions.

  \bigskip

  \textbf{Why it matters for brain mechanics:}
  \begin{itemize}
    \item The warp field $\bfu(\bfx)$ from diffeomorphic registration
      approximates the \emph{empirical tissue displacement}
    \item Jacobian $J(\bfx) = \det(\bfI + \nabla\bfu)$ encodes
      \emph{local volume change}
    \item These provide rare \alert{validation data} for FE model
      predictions
  \end{itemize}

  \bigskip

  \textbf{Challenge:} The resection cavity is a topological change ---
  tissue that existed pre-op is gone post-op. Standard registration
  tries to ``explain'' the cavity, producing artifacts.

\end{frame}

% ---- Frame 2 ----
\begin{frame}{The Cavity Masking Problem}

  \begin{columns}[T]
    \column{0.5\textwidth}
    \textbf{Without masking:}
    \begin{itemize}
      \item SyN tries to deform healthy tissue into the surgical void
      \item Produces unrealistic deformation near cavity
      \item Contaminates displacement and Jacobian maps
    \end{itemize}

    \bigskip

    \textbf{With masking:}
    \begin{itemize}
      \item Cavity region excluded from cost function
      \item SyN optimizes only on intact tissue
      \item Biologically plausible deformation field
    \end{itemize}

    \column{0.45\textwidth}
    % FIGURE SLOT: cavity mask overlay
    \begin{center}
      \fbox{\parbox{0.9\textwidth}{\centering\small
        \textit{GSlides:} Cavity mask\\
        overlay on post-op T1w\\[2pt]
        \texttt{FV-10\_cavity\_mask.png}
      }}
    \end{center}
  \end{columns}

\end{frame}

% ---- Frame 3 ----
\begin{frame}{Automatic Cavity Detection: Three Tiers}

  \textbf{Tier 1 (default):} Tumor-mask prior
  \begin{itemize}
    \item Resample the pre-op tumor mask into post-op T1w space
    \item Threshold at $p > 0.5$ $\to$ binary cavity seed
    \item Reliable when a pre-op tumor mask exists
  \end{itemize}

  \medskip

  \textbf{Tier 2 (fallback):} SynthSeg segmentation difference
  \begin{itemize}
    \item Compare pre-op and post-op \texttt{aparc+aseg}: brain $\to$
      CSF = cavity
    \item Sensitive to grid-size and label-convention mismatches
  \end{itemize}

  \medskip

  \textbf{Tier 3 (optional):} RESSEG deep learning model
  \cite{PerezGarcia2021}
  \begin{itemize}
    \item Enabled via \texttt{--resseg} flag
    \item Trained for postoperative cavity segmentation, but may
      mislocalize on atypical anatomies
  \end{itemize}

  \medskip

  All tiers produce a binary mask $\to$ inverted for ANTs \texttt{-x}
  flag (1 = healthy, 0 = cavity).
  Method-aware caching: switching tier auto-clears the mask
  \emph{and} all downstream outputs.

\end{frame}

% ---- Frame 3b: Diagnostic case study ----
\begin{frame}{Case Study: When Cavity Masking Fails}

  On sub-PAT16, RESSEG (Tier~3) produced a mask \alert{$\sim$50\,mm from
  the actual resection cavity}:

  \medskip

  \begin{tabular}{lrr}
    \toprule
    \textbf{Mask} & \textbf{Centre of mass (mm)} & \textbf{Volume} \\
    \midrule
    Pre-op tumor (Aerts) & $(30,\; 37,\; {-}13)$ & 45.4\,cm$^3$ \\
    RESSEG cavity        & $(53,\; {-}5,\; {-}21)$ & 18.6\,cm$^3$ \\
    Tumor-prior (Tier~1) & $(31,\; 36,\; {-}13)$ & 45.2\,cm$^3$ \\
    \bottomrule
  \end{tabular}

  \medskip

  \textbf{Consequence cascade:}
  \begin{itemize}
    \item Overlap between RESSEG mask and warped tumor: \alert{0\%}
    \item SyN excluded normal tissue, tried to match the actual cavity
    \item Global T1w correlation: \alert{0.66} (expected $> 0.85$)
    \item All 21 longitudinal outputs unreliable
  \end{itemize}

  \medskip

  \alert{Lesson:} Empirical verification of cavity masks is essential
  before downstream mechanics analysis.

\end{frame}

% ---- Frame 4 ----
\begin{frame}[fragile]{ANTs SyN with Cost Function Masking}

  \begin{lstlisting}
antsRegistrationSyN.sh \
  -d 3 \
  -f postop_T1w.nii.gz \     # fixed (post-op)
  -m preop_T1w.nii.gz  \     # moving (pre-op)
  -x cavity_mask_inv.nii.gz \ # cost function mask
  -t s \                      # SyN transform
  -n 8                        # threads
  \end{lstlisting}

  \bigskip

  \textbf{ANTs SyN} \cite{Avants2008}: symmetric diffeomorphic
  registration with cross-correlation.

  \bigskip

  \textbf{Outputs:}
  \begin{itemize}
    \item Forward warp: $\bfu(\bfx)$ (pre-op $\to$ post-op)
    \item Inverse warp: $\bfu^{-1}(\bfx)$ (post-op $\to$ pre-op)
    \item Affine component: initial linear alignment
  \end{itemize}

\end{frame}

% ---- Frame 5 ----
\begin{frame}{Deformation Field Outputs}

  From the SyN warp field we compute:

  \bigskip

  \begin{tabular}{lll}
    \toprule
    \textbf{Output} & \textbf{Definition} & \textbf{Interpretation} \\
    \midrule
    Displacement $|\bfu(\bfx)|$ &
      $\sqrt{u_1^2 + u_2^2 + u_3^2}$ &
      Tissue motion (mm) \\[4pt]
    Jacobian $J(\bfx)$ &
      $\det(\bfI + \nabla\bfu)$ &
      $J > 1$: expansion \\
    & & $J < 1$: contraction \\[4pt]
    Cavity mask & Binary &
      Resection void \\
    \bottomrule
  \end{tabular}

  \bigskip

  \begin{itemize}
    \item Displacement: where did tissue move, and how far?
    \item Jacobian: which regions expanded (e.g., ventricles) or
      contracted (e.g., near resection)?
    \item ROI-level statistics: mean displacement and Jacobian
      per brain region
  \end{itemize}

\end{frame}

% ---- Frame 6 ----
\begin{frame}{Warping Pre-op Derivatives to Post-op Space}

  With the SyN warp, we can bring \emph{any} pre-op derivative into
  post-op space for consistent voxel-wise comparison:

  \bigskip

  \begin{itemize}
    \item T1w (linear interpolation)
    \item FA, MD, AD, RD maps (linear interpolation)
    \item \texttt{aparc+aseg} (nearest-neighbor interpolation)
  \end{itemize}

  \bigskip

  \textbf{Result:} $\Delta\text{FA}(\bfx) = \text{FA}_{\text{post}}(\bfx)
  - \text{FA}_{\text{pre}\to\text{post}}(\bfx)$ is now a geometrically
  correct voxel-wise comparison, not contaminated by brain shift.

  \bigskip

  % FIGURE SLOT: warped pre-op T1w vs post-op T1w (checkerboard)
  \begin{center}
    \fbox{\parbox{0.7\textwidth}{\centering\small
      \textit{GSlides:} Warped pre-op T1w aligned with post-op T1w\\
      + displacement magnitude overlay\\[2pt]
      \texttt{FV-11\_warped\_alignment.png}
    }}
  \end{center}

\end{frame}

% ---- Frame 7 ----
\begin{frame}{Brain Mechanics Relevance of Deformation Fields}

  \begin{itemize}
    \item \textbf{Displacement field} $\bfu(\bfx)$: empirical tissue
      motion between pre-op and post-op
      \begin{itemize}
        \item Validation target for FE model predictions
        \item Spatial pattern reveals deformation mode
          (compression, shear, rotation)
      \end{itemize}
    \item \textbf{Jacobian} $J(\bfx)$: local volume change
      \begin{itemize}
        \item $J > 1$ near ventricles $\Rightarrow$ post-op expansion
        \item $J < 1$ near resection $\Rightarrow$ tissue collapse
        \item Regions where $J \neq 1$ require large-deformation
          formulations
      \end{itemize}
    \item \textbf{Cavity mask}: defines a new boundary condition surface
      for the post-op FE model
  \end{itemize}

  \bigskip

  \alert{These outputs bridge imaging and mechanics:}\\
  registration $\to$ displacement $\to$ strain $\to$ stress
  (via constitutive law).

\end{frame}
